% Appendix source input for exdqlm-jss.tex.
% This file is intentionally not a standalone JSS manuscript.

\section{Technical overview and notation}
\label{sec:app_scope}

The appendix aims to connect the package interface discussed in the article with the
posterior targets and computations implemented in \pkg{exdqlm}. It records the
distributional conventions, compact dynamic and static posterior targets, the
Laplace--delta variational Bayes (LDVB) scale--skewness block, the static
\code{rhs\_ns} shrinkage block, diagnostic and scoring quantities, and the
posterior-predictive synthesis algorithm. Standard Gaussian filtering,
smoothing, and backward-sampling recursions are cited rather than rederived.
Function-level arguments and backend controls are documented in the package
help pages, and the replication materials provide executable code for the
reported figures, tables, and benchmark quantities.

\subsection{Distributional conventions}
\label{sec:app_distributional_conventions}

Let $p_0\in(0,1)$ be the target quantile. The normal distribution
$\N(m,V)$ is parameterized by its mean and variance. The inverse-gamma
distribution uses the shape-rate convention
\begin{equation}
x\sim\IG(a,b)
\quad\Longleftrightarrow\quad
p(x)=\frac{b^a}{\Gamma(a)}x^{-a-1}\exp(-b/x),\qquad x>0.
\end{equation}
The exponential distribution $\Exp(\text{rate}=1/\sigma)$ has mean $\sigma$.
The notation $\TNp(m,V)$ denotes $\N(m,V)$ truncated to $(0,\infty)$.

The asymmetric Laplace (AL) mixture representation used by the package is
\begin{align}
v &\mid \sigma \sim \Exp(\text{rate}=1/\sigma),\\
y \mid \eta,\sigma,v
&\sim \N\{ \eta + A(p_0)v,\ \sigma B(p_0)v \},
\end{align}
where $\eta$ is the target-quantile location and
\begin{equation}
A(p)=\frac{1-2p}{p(1-p)}, \qquad B(p)=\frac{2}{p(1-p)}.
\label{eq:al_AB}
\end{equation}

The extended asymmetric Laplace (exAL) representation adds a latent
$s\sim\TNp(0,1)$. For $\gamma\in(L,U)$, define
\begin{align}
g(\gamma) &= 2\Phi(-|\gamma|)\exp(\gamma^2/2),\\
p_\gamma &= \Ind(\gamma<0) +
  \frac{p_0-\Ind(\gamma<0)}{g(\gamma)},\\
A_\gamma &= A(p_\gamma),\qquad
B_\gamma = B(p_\gamma),\qquad
C_\gamma = \{\Ind(\gamma>0)-p_\gamma\}^{-1}.
\label{eq:exal_ABC}
\end{align}
The bounds $L<0<U$ are the roots $g(L)=1-p_0$ and $g(U)=p_0$. Conditional on
$(v,s,\sigma,\gamma)$,
\begin{equation}
y\mid \eta,\sigma,\gamma,v,s
\sim
\N\left(\eta + C_\gamma\sigma|\gamma|s + A_\gamma v,\
\sigma B_\gamma v \right).
\label{eq:exal_aug}
\end{equation}
The package uses $\sigma\sim\IG(a_\sigma,b_\sigma)$ and a proper truncated
Student-$t$ prior $\pigamma(\gamma)$ for $\gamma$ on $(L,U)$:
\begin{equation}
\pigamma(\gamma)
\propto
s_\gamma^{-1}t_{\nu_\gamma}\{(\gamma-m_\gamma)/s_\gamma\}
\Ind(L<\gamma<U),
\label{eq:gamma_prior}
\end{equation}
with the normalizing constant determined by the Student-$t$ distribution
function over $(L,U)$.

The generalized inverse Gaussian (GIG) distribution is parameterized as
\begin{equation}
x\sim\GIG(\lambda,\chi,\psi_{\mathrm{GIG}})
\quad\Longleftrightarrow\quad
p(x)\propto x^{\lambda-1}
\exp\left\{-\frac{1}{2}\left(\frac{\chi}{x}
+\psi_{\mathrm{GIG}} x\right)\right\},\qquad x>0.
\label{eq:gig_parameterization}
\end{equation}
For $r\in\R$,
\begin{equation}
\E(x^r)=
\left(\frac{\chi}{\psi_{\mathrm{GIG}}}\right)^{r/2}
\frac{K_{\lambda+r}(\sqrt{\chi\psi_{\mathrm{GIG}}})}
{K_{\lambda}(\sqrt{\chi\psi_{\mathrm{GIG}}})},
\label{eq:gig_moments}
\end{equation}
where $K_\nu(\cdot)$ is the modified Bessel function of the second kind.

\section{Posterior targets and model blocks}

\subsection{Dynamic DQLM under the AL likelihood}
\label{sec:dyn_dqlm}

For the dynamic quantile linear model (DQLM), the observation and evolution
equations are
\begin{align}
y_t\mid\vect{\theta}_t,\sigma,v_t
&\sim
\N\left(\mat{F}_t^\top\vect{\theta}_t+A(p_0)v_t,\
\sigma B(p_0)v_t\right),\\
\vect{\theta}_t\mid\vect{\theta}_{t-1}
&\sim \N(\mat{G}_t\vect{\theta}_{t-1},\mat{W}_t),\\
v_t\mid\sigma &\sim \Exp(\text{rate}=1/\sigma).
\label{eq:dyn_dqlm_hierarchy}
\end{align}
The posterior target is proportional to
\begin{equation}
p(\vect{\theta}_{0:T},\vect{v},\sigma\mid\vect{y})
\propto
p(\vect{\theta}_0)p(\sigma)
\prod_{t=1}^T
p(y_t\mid\vect{\theta}_t,\sigma,v_t)
p(v_t\mid\sigma)
p(\vect{\theta}_t\mid\vect{\theta}_{t-1}).
\label{eq:dyn_dqlm_joint}
\end{equation}
Conditioning on $(\sigma,\vect{v})$, the state sequence is a Gaussian
state-space model with observation offset $A(p_0)v_t$ and variance
$\sigma B(p_0)v_t$. The package uses standard Kalman filtering,
Rauch--Tung--Striebel smoothing, and forward-filtering backward-sampling
recursions for this Gaussian block \citep{west1997bayesian}.
\label{sec:gaussian_state_update}

\subsection{Dynamic exDQLM under the exAL likelihood}
\label{sec:dyn_exdqlm}

The dynamic exDQLM augments the AL hierarchy with $s_t\sim\TNp(0,1)$ and
the exAL constants in \eqref{eq:exal_ABC}:
\begin{align}
y_t\mid\vect{\theta}_t,\sigma,\gamma,v_t,s_t
&\sim
\N\left(
\mat{F}_t^\top\vect{\theta}_t
+C_\gamma\sigma|\gamma|s_t+A_\gamma v_t,\
\sigma B_\gamma v_t
\right),\\
\vect{\theta}_t\mid\vect{\theta}_{t-1}
&\sim \N(\mat{G}_t\vect{\theta}_{t-1},\mat{W}_t),\\
v_t\mid\sigma&\sim\Exp(\text{rate}=1/\sigma),\qquad
s_t\sim\TNp(0,1).
\label{eq:dyn_exdqlm_obs}
\end{align}
The posterior target is
\begin{align}
p(\vect{\theta}_{0:T},\vect{v},\vect{s},\sigma,\gamma\mid\vect{y})
&\propto
p(\vect{\theta}_0)p(\sigma)\pigamma(\gamma)
\prod_{t=1}^T p(y_t\mid\vect{\theta}_t,\sigma,\gamma,v_t,s_t)
\nonumber\\
&\quad{}\times
\prod_{t=1}^T
p(v_t\mid\sigma)p(s_t)
p(\vect{\theta}_t\mid\vect{\theta}_{t-1}).
\label{eq:dyn_exdqlm_joint}
\end{align}
Given $(\sigma,\gamma,\vect{v},\vect{s})$, the state sequence is again a
Gaussian state-space model, now with offset
$A_\gamma v_t+C_\gamma\sigma|\gamma|s_t$ and variance
$\sigma B_\gamma v_t$.

\subsection{Transfer-function state augmentation}
\label{sec:transfer_function}

For an input $x_t$ and transfer lag depth $r$, the transfer-function block
adds states
\[
\vect{\psi}_t=(\psi_{0,t},\ldots,\psi_{r,t})^\top,\qquad
\vect{\zeta}_t=(\zeta_{0,t},\ldots,\zeta_{r,t})^\top,
\]
and evolves the accumulated input effect with a lag-decay rate $\lambda$.
In the package, $\lambda$ is user-specified or selected externally by a grid
search. The current formulation uses a common lag-decay rate within a
transfer block. The state augmentation used by
\code{exdqlmTransferMCMC()} and \code{exdqlmTransferLDVB()} is
\begin{equation}
\vect{\theta}^{\mathrm{aug}}_t
=
(\vect{\theta}^{\mathrm{base}\top}_t,
\vect{\zeta}^{\top}_t,\vect{\psi}^{\top}_t)^\top,
\label{eq:tf_aug_state}
\end{equation}
with the base state-space matrices expanded to include the transfer states.
After augmentation, the fitting algorithms are the same dynamic MCMC or LDVB
routines used for the corresponding exDQLM.

\subsection{Static AL and exAL regression}
\label{sec:static_ridge}

Let $\mat{X}\in\R^{n\times p}$ have rows $\vect{x}_i^\top$ and regression
coefficient vector $\vect{\beta}$. Under a Gaussian/ridge prior,
\begin{equation}
\vect{\beta}\sim\N(\vect{b}_{\beta,0},\mat{V}_{\beta,0}).
\label{eq:static_beta_ridge}
\end{equation}
The static exAL likelihood is
\begin{align}
y_i\mid\vect{\beta},\sigma,\gamma,v_i,s_i
&\sim
\N\left(
\vect{x}_i^\top\vect{\beta}
+C_\gamma\sigma|\gamma|s_i+A_\gamma v_i,\
\sigma B_\gamma v_i
\right),\\
v_i\mid\sigma&\sim \Exp(\text{rate}=1/\sigma),\qquad
s_i\sim\TNp(0,1).
\label{eq:static_exal_hierarchy}
\end{align}
The posterior target is
\begin{align}
p(\vect{\beta},\sigma,\gamma,\vect{v},\vect{s}\mid\vect{y})
&\propto
p(\vect{\beta})p(\sigma)\pigamma(\gamma)
\prod_{i=1}^n
p(y_i\mid\vect{\beta},\sigma,\gamma,v_i,s_i)
p(v_i\mid\sigma)p(s_i).
\label{eq:static_exal_target}
\end{align}
The AL static regression is recovered by fixing $\gamma=0$ and dropping the
$s_i$ block. For fixed latent variables and scale parameters, the
coefficient update is the weighted Gaussian regression update; for
\code{beta_prior = "rhs\_ns"}, the Gaussian prior precision is replaced by
the expected \code{rhs\_ns} precision described in
Section~\ref{sec:static_rhs}.

\section{Laplace--delta VB implementation for exAL models}
\label{sec:ldvb_details}

The nonconjugate part of the exAL posterior is the joint scale--skewness
block $(\sigma,\gamma)$. The LDVB routines in \pkg{exdqlm} specialize the
Laplace--delta strategy of \citet{wang2013nonconjugatevb} to this block.
The dynamic variational family is
\begin{equation}
q(\vect{\theta}_{0:T},\vect{v},\vect{s},\sigma,\gamma)
=q(\vect{\theta}_{0:T})\prod_{t=1}^Tq(v_t)\prod_{t=1}^Tq(s_t)\,
q(\sigma,\gamma).
\label{eq:dyn_exdqlm_mf}
\end{equation}
The static family replaces $q(\vect{\theta}_{0:T})$ by
$q(\vect{\beta})$:
\begin{equation}
q(\vect{\beta},\vect{v},\vect{s},\sigma,\gamma)
=q(\vect{\beta})\prod_{i=1}^nq(v_i)\prod_{i=1}^nq(s_i)\,
q(\sigma,\gamma).
\label{eq:static_exal_mf}
\end{equation}

The conjugate variational factors require moments of functions of
$(\sigma,\gamma)$. In the dynamic case, define
\begin{align}
\kappa_1 &= \E\{1/(\sigma B_\gamma)\},&
\kappa_2 &= \E\{A_\gamma/(\sigma B_\gamma)\},&
\kappa_3 &= \E\{A_\gamma^2/(\sigma B_\gamma)\},\\
\kappa_4 &= \E\{C_\gamma|\gamma|/B_\gamma\},&
\kappa_5 &= \E\{\sigma C_\gamma^2\gamma^2/B_\gamma\},&
\kappa_6 &= \E\{A_\gamma C_\gamma|\gamma|/B_\gamma\}.
\label{eq:dyn_exdqlm_kappa}
\end{align}
For the dynamic state factor, let $m_{1/v,t}=\E(v_t^{-1})$ and
$\bar s_t=\E(s_t)$. The Gaussian pseudo-observation has precision
\begin{equation}
\phi_t=\kappa_1m_{1/v,t}
\label{eq:dyn_exdqlm_vb_precision}
\end{equation}
and pseudo-response
\begin{equation}
y_t^{\mathrm{VB}}
=
\frac{
y_t\phi_t-\kappa_2-\kappa_4\bar s_t m_{1/v,t}
}{\phi_t}.
\label{eq:dyn_exdqlm_vb_pseudo}
\end{equation}
Thus the state factor is updated by the same Gaussian state-space recursions
as in Section~\ref{sec:dyn_dqlm}.

For static regression, the coefficient factor is Gaussian. With row vectors
$\vect{x}_i^\top$,
\begin{align}
\omega_i&=\kappa_1\E(v_i^{-1}),\\
\eta_i&=
y_i\kappa_1\E(v_i^{-1})
-\kappa_4\E(v_i^{-1})\E(s_i)
-\kappa_2,\\
\mat{\Sigma}_\beta
&=\left(\mat{V}_{\beta,0}^{-1}
+\mat{X}^\top\diag(\omega_1,\ldots,\omega_n)\mat{X}\right)^{-1},\\
\vect{m}_\beta
&=\mat{\Sigma}_\beta
\left(\mat{V}_{\beta,0}^{-1}\vect{b}_{\beta,0}
+\mat{X}^\top\vect{\eta}\right),
\label{eq:static_exal_vb_beta}
\end{align}
where $\vect{\eta}=(\eta_1,\ldots,\eta_n)^\top$.

The latent-variable factors remain closed form conditional on the
Laplace--delta moments. If $m_{\eta,t}=\E(\mat{F}_t^\top\vect{\theta}_t)$,
$S_{\eta,t}=\Var(\mat{F}_t^\top\vect{\theta}_t)$, and
$S_{\Delta,t}=(y_t-m_{\eta,t})^2+S_{\eta,t}$, then the dynamic updates are
\begin{align}
q(v_t)&=\GIG\left(
\frac{1}{2},\
\kappa_1S_{\Delta,t}-2\kappa_4(y_t-m_{\eta,t})\E(s_t)
+\kappa_5\E(s_t^2),\
\kappa_3+2\E(\sigma^{-1})
\right),\\
q(s_t)&=\TNp(\mu_{s,t}^{\mathrm{VB}},V_{s,t}^{\mathrm{VB}}),\\
V_{s,t}^{\mathrm{VB}}
&=\left(1+\kappa_5\E(v_t^{-1})\right)^{-1},\\
\mu_{s,t}^{\mathrm{VB}}
&=V_{s,t}^{\mathrm{VB}}
\left\{\kappa_4(y_t-m_{\eta,t})\E(v_t^{-1})-\kappa_6\right\}.
\label{eq:dyn_exdqlm_vb_vs}
\end{align}
The static updates replace $m_{\eta,t}$ and $S_{\eta,t}$ by
$\vect{x}_i^\top\vect{m}_\beta$ and
$\vect{x}_i^\top\mat{\Sigma}_\beta\vect{x}_i$.

The nonconjugate factor is approximated on transformed coordinates
\begin{equation}
\ell_\sigma=\log\sigma,\qquad
z_\gamma=\logit\left(\frac{\gamma-L}{U-L}\right).
\label{eq:dyn_exdqlm_ld_transform}
\end{equation}
If $r=\expit(z_\gamma)$, the transformed objective is
\begin{align}
h_z(\ell_\sigma,z_\gamma)
&=
\E_{-(\sigma,\gamma)}
\{\log p(\vect{y},\text{latent variables},\text{parameters})\}
\nonumber\\
&\quad
+\ell_\sigma+\log(U-L)+\log r+\log(1-r).
\label{eq:dyn_exdqlm_ld_objective}
\end{align}
The package maximizes this objective, uses the local Hessian to form a
Gaussian approximation on $(\ell_\sigma,z_\gamma)$, and evaluates needed
expectations by the delta method. The returned LDVB objects include the
transformed mode, approximate covariance, derived moments, ELBO trace, and
approximate posterior draws when requested.

The dynamic and static exAL ELBOs have the same block structure:
\begin{equation}
\elbo(q)=
\elbo_{\mathrm{like}}+\elbo_{\mathrm{Gaussian}}
+\elbo_v+\elbo_s+\elbo_{\sigma\gamma}
+H(q_{\mathrm{Gaussian}})
+\sum H\{q(v)\}+\sum H\{q(s)\}+H\{q(\sigma,\gamma)\}.
\label{eq:dyn_exdqlm_elbo}
\end{equation}
For dynamic models, the Gaussian block is the state trajectory
$q(\vect{\theta}_{0:T})$; for static models, it is $q(\vect{\beta})$.
The entropy terms use the distributional conventions stated in
Section~\ref{sec:app_distributional_conventions}. For AL/DQLM fits,
$\gamma=0$, the $s$ block is omitted, and $q(\sigma)$ is inverse-gamma.

\section{Static regression with the Nishimura--Suchard regularized horseshoe}
\label{sec:static_rhs}

The \code{rhs\_ns} option implements the conditionally conjugate
regularized horseshoe formulation of \citet{nishimura2023shrunken} for
static regression. Let $\mathcal{A}\subseteq\{1,\ldots,p\}$ be the set of
coefficients subject to shrinkage; by default an intercept is not shrunk. For
$j\in\mathcal{A}$,
\begin{align}
\beta_j\mid\tau^2,\lambda_j^2,\zeta^2
&\propto
\N(\beta_j\mid0,\tau^2\lambda_j^2)
\N(\beta_j\mid0,\zeta^2),\\
\lambda_j^2\mid\nu_j &\sim \IG(1/2,1/\nu_j),&
\nu_j&\sim\IG(1/2,1),\\
\tau^2\mid\xi &\sim\IG(1/2,1/\xi),&
\xi&\sim\IG(1/2,1/\tau_0^2),\\
\zeta^2&\sim\IG(a_\zeta,b_\zeta),
\label{eq:rhs_ns_hierarchy}
\end{align}
unless the slab variance $\zeta^2$ is fixed by the user. The active prior
precision contribution is
\begin{equation}
\omega^{\mathrm{prior}}_j
=\frac{1}{\tau^2\lambda_j^2}+\frac{1}{\zeta^2}.
\label{eq:rhs_ns_precision}
\end{equation}
The static Gaussian coefficient update uses this precision, or its
variational expectation, on the active diagonal entries.

Conditioning on $\vect{\beta}$, the shrinkage MCMC block has closed-form
updates:
\begin{align}
\lambda_j^2\mid\cdot
&\sim
\IG\left(1,\ \frac{1}{\nu_j}+\frac{\beta_j^2}{2\tau^2}\right),\\
\nu_j\mid\cdot
&\sim
\IG\left(1,\ 1+\frac{1}{\lambda_j^2}\right),\\
\tau^2\mid\cdot
&\sim
\IG\left(\frac{|\mathcal{A}|+1}{2},\
\frac{1}{\xi}
+\frac{1}{2}\sum_{j\in\mathcal{A}}\frac{\beta_j^2}{\lambda_j^2}
\right),\\
\xi\mid\cdot
&\sim
\IG\left(1,\ \frac{1}{\tau_0^2}+\frac{1}{\tau^2}\right),\\
\zeta^2\mid\cdot
&\sim
\IG\left(a_\zeta+\frac{|\mathcal{A}|}{2},\
b_\zeta+\frac{1}{2}\sum_{j\in\mathcal{A}}\beta_j^2
\right),
\label{eq:rhs_ns_mcmc}
\end{align}
with the final line omitted when $\zeta^2$ is fixed.

The mean-field VB block has inverse-gamma factors
\[
q(\lambda_j^2),\quad q(\nu_j),\quad q(\tau^2),\quad q(\xi),
\quad q(\zeta^2)\ \text{if }\zeta^2\text{ is random}.
\]
Let $\bar\beta_j^2=\E_q(\beta_j^2)$. The rate updates are
\begin{align}
b_{\lambda j}
&=\frac{1}{2}\bar\beta_j^2\,\E\{(\tau^2)^{-1}\}
+\E(\nu_j^{-1}),\\
b_{\nu j}
&=1+\E\{(\lambda_j^2)^{-1}\},\\
b_\tau
&=\frac{1}{2}\sum_{j\in\mathcal{A}}\bar\beta_j^2
\E\{(\lambda_j^2)^{-1}\}+\E(\xi^{-1}),\\
b_\xi
&=\frac{1}{\tau_0^2}+\E\{(\tau^2)^{-1}\},\\
a_\zeta^\star
&=a_\zeta+\frac{|\mathcal{A}|}{2},\qquad
b_\zeta^\star=b_\zeta+\frac{1}{2}\sum_{j\in\mathcal{A}}\bar\beta_j^2.
\label{eq:rhs_ns_vb_updates}
\end{align}
The shape parameters are
$a_{\lambda j}=a_{\nu j}=a_\xi=1$ and
$a_\tau=(|\mathcal{A}|+1)/2$.

The \code{rhs\_ns} ELBO is the sum of coefficient-prior, scale-prior,
slab-prior, and inverse-gamma entropy terms:
\begin{equation}
\elbo_{\mathrm{rhs\_ns}}
=
\elbo_{\beta,\mathrm{hs}}
+\elbo_{\beta,\mathrm{slab}}
+\elbo_{\lambda\mid\nu}
+\elbo_{\nu}
+\elbo_{\tau\mid\xi}
+\elbo_{\xi}
+\elbo_{\zeta}
+\elbo_{\mathrm{ent}}
+\elbo_{\beta_0}.
\label{eq:rhs_ns_elbo}
\end{equation}
All terms are available in closed form because each scale factor is
inverse-gamma. This block is added to the static AL or exAL ELBO, whose
likelihood and scale--skewness components are handled by the surrounding
static fitting routine. The fitted static diagnostic object uses the resulting
posterior coefficient draws to produce the coefficient-level summaries shown
in Example 4.

\section{Forecasting, diagnostics, and posterior-predictive synthesis}
\label{sec:forecast_diagnostics_synthesis}

\subsection{Forecasting}

At forecast origin $\tilde t$, the dynamic state forecast recursion gives
\[
\vect{\theta}_{\tilde t+k}^{(m)}
\sim
\N\{a_{\tilde t}^{(m)}(k),R_{\tilde t}^{(m)}(k)\},
\qquad k=1,\ldots,K,
\]
for each posterior or variational draw $m$. Posterior predictive draws are
then generated from
\begin{equation}
Y_{\tilde t+k}^{\mathrm{rep},(m)}
\sim
\exAL_{p_0}\{
\mat{F}_{\tilde t+k}^\top\vect{\theta}_{\tilde t+k}^{(m)},
\sigma^{(m)},\gamma^{(m)}
\}.
\label{eq:forecast_postpred}
\end{equation}
For the AL/DQLM special case, $\gamma^{(m)}=0$. The function
\code{exdqlmForecast()} implements this recursion, and the standard
\code{predict()} method for dynamic fitted objects calls the same forecast
constructor.

\subsection{Diagnostics}

The fitted dynamic diagnostic sequence is based on maximum a posteriori (MAP)
one-step-ahead standardized forecast errors. If $\widehat f_t$ and
$\widehat q_t$ are the MAP one-step-ahead location and variance, then
\begin{equation}
\widehat e_t = \frac{y_t-\widehat f_t}{\sqrt{\widehat q_t}},
\qquad
\widehat u_t = \Phi(\widehat e_t).
\label{eq:pit_plugin}
\end{equation}
The diagnostic plot method uses $\{\widehat e_t\}$ and $\{\widehat u_t\}$ for
the QQ plot, autocorrelation function plot, and standardized forecast-error
plot. The reported Kullback--Leibler (KL) values are deterministic
one-dimensional calibration diagnostics comparing the empirical distribution
of $\{\widehat e_t\}$ with $N(0,1)$; the package uses a deterministic normal
quantile grid rather than a random reference sample.

For posterior predictive draws $y_t^{(1)},\ldots,y_t^{(M)}$, let
$\widehat q_t(\tau_k)$ be the empirical posterior predictive quantile at
level $\tau_k$. The package approximates the continuous ranked probability
score (CRPS) by the integrated quantile-score identity
\begin{equation}
\widehat{\CRPS}_t
=
2\sum_{k=1}^{K}
w_k
\rho_{\tau_k}\{y_t^{\mathrm{obs}}-\widehat q_t(\tau_k)\},
\label{eq:crps_estimator}
\end{equation}
where the weights sum to one \citep{gneiting2007,laio2007verification}. The
posterior predictive loss criterion (PPLC) at the fitted quantile level is
\begin{equation}
\widehat{\PPLC}
=
\sum_t
\frac{1}{M}\sum_{m=1}^M
\rho_{p_0}\{y_t^{\mathrm{obs}}-y_t^{\mathrm{rep},(m)}\}.
\label{eq:pplc_estimator}
\end{equation}
The held-out forecast diagnostic function applies the same CRPS calculation
to forecast objects returned with posterior predictive draws.

\subsection{Posterior-predictive synthesis}

The function \code{quantileSynthesis()} combines posterior predictive draws
from separately fitted models at ordered levels $p_1<\cdots<p_K$. The result
is a post hoc predictive distribution, not a joint multi-quantile posterior.

\begin{appalgorithm}{Posterior-predictive synthesis}
\label{alg:predictive-synthesis}
\par\noindent\emph{Input:} posterior predictive draw matrices or fitted
forecast objects from models at ordered quantile levels
$p_1<\cdots<p_K$.
\par\noindent\emph{Output:} synthesized posterior predictive draws and
summary intervals on a common predictive scale.
\begin{enumerate}[label=(\arabic*)]
\item Extract posterior predictive draw matrices for each fitted quantile.
\item For each time point and draw index, form the quantile-specific
predictive curve over $p_1,\ldots,p_K$.
\item Apply the requested monotonicity correction when fitted predictive
quantiles cross.
\item Interpolate across quantile levels to obtain draws from one predictive
distribution.
\item Return synthesized posterior predictive draws and summary intervals.
\end{enumerate}
\end{appalgorithm}

\section{Laplace--delta and backend notes}
\label{sec:numerical_blocks}

Let $z$ denote the transformed nonconjugate parameter vector with mode
$\hat z$ and negative Hessian $\mat{H}$ at the mode. The LDVB block uses
\begin{equation}
q(z)=\N(\hat z,\mat{H}^{-1}).
\label{eq:ld_gaussian}
\end{equation}
For a smooth scalar function $h(z)$,
\begin{equation}
\E\{h(z)\}
\approx
h(\hat z)
+\frac{1}{2}\tr\{\nabla^2h(\hat z)\mat{H}^{-1}\}.
\label{eq:delta_generic}
\end{equation}
The entropy contribution for the transformed Gaussian block is
\begin{equation}
H\{q(z)\}
=
\frac{1}{2}\log\{(2\pi e)^d|\mat{H}^{-1}|\}.
\label{eq:entropy_ld}
\end{equation}
The Jacobian terms for transforming from $(\ell_\sigma,z_\gamma)$ back to
$(\sigma,\gamma)$ are included in the objective in
\eqref{eq:dyn_exdqlm_ld_objective}. Other ELBO terms use standard entropy
identities for Gaussian, inverse-gamma, GIG, and positive truncated-normal
factors under the parameterizations in
Section~\ref{sec:app_distributional_conventions}. Those bookkeeping terms are
implemented in the package and exercised by the package tests rather than
rederived here.

\label{sec:slice}

The MCMC routines use bounded one-dimensional slice sampling for selected
nonconjugate skewness updates. This is a standard stepping-out/shrinkage
slice sampler restricted to the admissible interval $(L,U)$ for $\gamma$; the
target kernels are the conditional log posteriors induced by the exAL
augmentation.

\label{sec:app_backend_options}
The package also exposes backend options for C++ Kalman recursions, builders,
latent-variable samplers, posterior predictive simulation, and MCMC routing.
These include the \code{exdqlm.use\_cpp\_*} option family, such as
\code{exdqlm.use\_cpp\_kf}, and the \code{exdqlm.cpp\_threads} thread cap.
They are documented in the package help pages and README because they are
runtime controls rather than additional model assumptions.
